\chapter{Background and Related Work}
\label{ch:background}

This chapter places the thesis in five bodies of work: dialect and accent
identification (\cref{sec:bg-did}); pretrained speech representations
(\cref{sec:bg-ssl}); Bangla speech and dialect resources
(\cref{sec:bg-bangla}); leakage and shortcut learning in evaluation
(\cref{sec:bg-leakage}); and augmentation for speech (\cref{sec:bg-aug}).
\Cref{sec:bg-asr} covers the Bangla speech-recognition work on which the
lexical route builds, and \cref{sec:bg-gap} states the gap this thesis
addresses.

\section{Dialect and accent identification}
\label{sec:bg-did}

Dialect identification is usually framed alongside language identification
(LID) and speaker recognition. Modern systems in all three inherit the embedding
paradigm of x-vectors \citep{snyder2018xvectors}: a network trained to
discriminate classes produces a fixed-length utterance embedding on which a
light classifier operates. The kinship with speaker recognition matters for
this thesis. A representation good at separating speakers is also able to
separate dialect classes by recognising the speakers in them, which is why
speaker overlap between training and test data inflates dialect scores
(\cref{sec:bg-leakage}).

The largest DID benchmarks are for Arabic. The MGB-5 challenge's ADI17 task
asks systems to identify 17 Arabic dialects from broadcast speech
\citep{ali2019mgb5}, the closest analogue to the 9-way Bangla setting here.
\citet{khaled2025arabic} fine-tune a cross-lingual wav2vec~2.0 model on Arabic
dialect audio, the same design as this thesis's fine-tuned model, and
\citet{sullivan2023robustness} show that Arabic DID systems degrade sharply when
moved away from their training conditions. For another lower-resource language,
\citet{hamalainen2021finnish} compare audio and text for Finnish dialect
identification, a precedent for asking whether dialect lives in the sound or in
the words. At larger scale, VoxLingua107 \citep{valk2021voxlingua} builds a
107-language LID corpus from YouTube, illustrating both the power of found data
and its channel confounds, and \citet{kukk2022accented} study how accent
interacts with language identification---relevant to treating the standard
register as a class of its own.

\section{Pretrained speech representations}
\label{sec:bg-ssl}

Wav2vec~2.0 \citep{baevski2020wav2vec} learns speech representations from raw
audio by solving a contrastive task over masked, quantised latent features, and
is then fine-tuned for a downstream task. XLS-R \citep{babu2022xlsr} scales the
same objective to 436,000 hours of speech in 128 languages, Bangla included; its
300-million-parameter variant is the model fine-tuned in this thesis. Whisper
\citep{radford2023whisper} takes a different route: an encoder--decoder
Transformer trained with weak supervision on 680,000 hours of transcribed,
multilingual web audio. Its encoder consumes 80-channel log-mel spectrograms
through two convolutional layers and a Transformer stack; this thesis uses the
encoder of Whisper-small as a frozen feature extractor.

The SUPERB benchmark \citep{yang2021superb} showed that frozen self-supervised
representations with light task heads transfer well across many speech tasks,
which motivates the frozen-embedding baseline here. Probing work on accent
identification \citep{yang2023accent} shows that fine-tuning a wav2vec~2.0
model for accent reshapes how its layers encode prosody---external support for
this thesis's explanation of why pitch and tempo augmentation hurt dialect
identification (\cref{sec:disc-augment}). The pseudo-speaker clustering in
\cref{ch:method} uses ECAPA-TDNN speaker embeddings \citep{desplanques2020ecapa}.

\section{Bangla speech and dialect resources}
\label{sec:bg-bangla}

Bangla speech processing is widely described as low-resource despite the
language's size and dialectal diversity \citep{bengaliasrsurvey2021}. The
closest prior work to the first version of this project classifies Bengali
accents from hand-crafted acoustic features with classical and neural
classifiers \citep{badhon2020accent}, and \citet{alam2023accent} classify Bangla
accent variation from audio with deep neural networks. The corpus used in this
thesis was gathered for those two studies (\cref{ch:dataset}). Text-based Bangla dialect classification
has progressed further: on the BanglaDial dataset, transformer models such as
BanglaBERT reach a reported score of about 89\,\% when classifying written
dialect text from 11 regions \citep{bangladial2025}. That result is a precedent for the lexical
route, with one difference: its text is written by people, whereas the lexical
route in this thesis has to obtain its text from speech recognition.

Regional speech datasets are appearing: a five-region speech dataset
\citep{bdregionalspeech}, the BRADS and BRWDS audio-and-text collections
\citep{brads2025}, and the Ben-10 and RegSpeech12 corpora discussed in
\cref{sec:bg-asr}. Large-scale standard Bangla speech is available in the
crowdsourced Bengali Common Voice corpus \citep{alam2022commonvoice} and in
OOD-Speech \citep{rakib2023oodspeech}.

\section{Leakage, shortcuts and honest evaluation}
\label{sec:bg-leakage}

Deep networks exploit whatever regularity separates the training classes,
including regularities that are not the intended signal
\citep{geirhos2020shortcut}. In audio, recording channel, background and session
are such shortcuts. SpurAudio \citep{spuraudio2026} benchmarks shortcut learning
in audio classification directly, and EchoHawk \citep{echohawk2026} shows, for
acoustic drone detection, that clips sharing a recording session share their
channel and background, and that session-disjoint evaluation collapses inflated
scores---the same phenomenon this thesis measures for dialects. A survey of
audio-classification practice finds evaluation protocols frequently
under-specified \citep{leakage2025iciis}. Together these legitimise reporting the
drop from a clip-level to a group-disjoint split as a finding in itself.

\section{Augmentation for speech}
\label{sec:bg-aug}

Speed and tempo perturbation is a standard and beneficial augmentation in speech
recognition \citep{ko2015augmentation}, where the target transcript does not
depend on prosody. SpecAugment \citep{park2019specaugment} masks time and
frequency bands of the spectrogram and is the canonical feature-level method.
\Cref{sec:res-augment} shows that the analogous waveform-level transforms reduce
dialect-identification accuracy, and \cref{sec:disc-augment} argues why the
recognition intuition does not carry over.

\section{Bangla speech recognition and dialect speech}
\label{sec:bg-asr}

The lexical route depends on three recent resources.

\paragraph{OOD-Speech.} \citet{rakib2023oodspeech} released 1,178 hours of
crowd-read standard Bangla from 22,645 speakers, with a 23-hour test set drawn
from 17 out-of-distribution domains. It was the training data of the 2023
Bengali.AI speech-recognition competition \citep{bengaliai2023competition},
whose winning entry---a fine-tuned Whisper-medium released as
\hf{bengaliAI/tugstugi_bengaliai-asr_whisper-medium} and referred to here as
\emph{Tugstugi}---is one of this thesis's recognisers. Its training transcripts
supply the standard-Bangla vocabulary in \cref{sec:method-lexical}.

\paragraph{Ben-10.} \citet{dipto2025ben10} introduce Ben-10, 78 hours of mostly
spontaneous speech from 394 speakers in ten districts, transcribed in the
dialect as spoken, and benchmark seven speech-recognition systems on it. Word
error rates range from 1.13 for zero-shot Whisper-large-v3 and 1.03 for
Google's cloud recogniser, through 0.87 for Hishab's Conformer and 0.81 for
Tugstugi, to 0.70 for Tugstugi fine-tuned on Ben-10. In a manual analysis of
about 50 dialect words per district, the authors find that the
standard-trained Tugstugi model reproduces 1--9\,\% of dialect words, against
16--53\,\% after fine-tuning on Ben-10, and describe a full-scale analysis as beyond their scope.
The paper is a speech-recognition benchmark, not a dialect-identification
study: it trains no classifier and reports no accuracy. In this thesis its
models serve as transcription front-ends, and its manual dialect-recall
analysis is repeated automatically over its full validation set
(\cref{sec:res-asr}).

\paragraph{Parallel dialect text.} Vashantor \citep{faria2023vashantor}
provides 32,500 sentence pairs translating regional Bangla from Chittagong,
Noakhali, Sylhet, Barishal and Mymensingh into standard Bangla. Because each
pair contains both forms, a regional word can be defined as one that appears on
the dialect side but not on the standard side. RegSpeech12
\citep{hassan2025regspeech12}, from the same group as Ben-10, extends coverage to
twelve dialects, and BanglaDialecto \citep{bangladialecto2024} provides Noakhali
speech with dialect and standard transcriptions; neither is used here, but both
would extend word lists to more of this corpus's classes.

The other recognisers used are Hishab's FastConformer \citep{nandi2023hishab},
trained on about 18,000 hours of mostly broadcast-news Bangla, Whisper-large-v3
\citep{radford2023whisper}, and a public wav2vec~2.0 XLS-R model fine-tuned for
Bangla \citep{arijitx2022wav2vec2}, which stands in for the benchmark's
``Wav2Vec2 (SCB)'' system whose checkpoint is not public.

\section{The gap}
\label{sec:bg-gap}

Prior Bangla dialect work uses either hand-crafted acoustic features on small
corpora or human-written dialect text. None applies pretrained speech models
under an evaluation that controls speaker and channel leakage, and none asks how
much dialect information survives the step from speech to text. This thesis
addresses both: it compares pretrained and hand-crafted acoustic
representations under a measured, leakage-aware protocol, and it tests whether
the words that speech recognition recovers can identify the dialect, alone or in
combination with the sound.
